\documentclass[11pt]{article}
\usepackage[margin=0.85in]{geometry}
\usepackage[T1]{fontenc}
\usepackage{amsmath,amssymb,booktabs,longtable,microtype,xurl,hyperref}
\setlength{\emergencystretch}{3em}
\hypersetup{colorlinks=true,urlcolor=blue,linkcolor=blue}
\title{Faster PCBF/CLF Optimization: Retained Changes and Rejected Trust-Growth Trials}
\author{Pan Dynamics Collision Avoidance Research}
\date{October 2, 2026}
\begin{document}\maketitle

\section{Measured outcome}
The previously worst recorded controller frame falls from 2.949 s to 0.427 s
in a paired replay. The final complete campaign has a 0.032910 s median,
0.327008 s 95th percentile, and a 0.514561 s maximum after the process's
first controller call. That first cold call takes 1.263946 s. Every figure
covers the complete controller call. The requested 0.1 s maximum is not met.

The retained implementation has thirteen of fourteen collision-free nominal
recoveries, fourteen eventual recoveries, and no controller interruption.
This is the same aggregate collision-free count as the preceding controller,
but the failed case changes: the 15 m/s straight crossing now passes with
0.088047 m clearance, while the 15 m/s turning crossing has intersample
contact. Equal aggregate counts must not be described as absence of a safety
regression in every individual scenario.

The analytic primary certificate and native quadratic conic optimization are
retained. The third proposal was implemented and tested: the objective-resolution
stopping rule is retained, but the added accurate-step trust growth is removed
because it produced another collision in the curved-crossing case. The fastest
trial is therefore not the final controller. No collision-buffer tuning,
scenario-specific switching, additional controller, or second CLF is introduced.

\section{Retained algorithm changes}
\subsection{Zero-primary optimality without a numerical solve}
For $J_s=\sum_i\xi_i$ with $\xi_i\ge0$, a fully feasible zero-slack anchor
proves $0\le J_s^\star\le0$. The implementation checks the zero correction
against every assembled inequality, equality, bound and terminal cone using
strict floating-point comparisons. A small positive residual or primary value
is not rounded to zero. The primary stage remains in the record with
\texttt{numericalSolve=false}; the CLF stage always executes. A collision-containing
or terminal-infeasible anchor cannot use this certificate.

\subsection{Equivalent native quadratic objective}
With $M$ holds and input correction scale $D=\operatorname{diag}(r_\delta,r_b)$,
\[
 R(\Delta U)=\frac{1}{2M}\sum_{i=0}^{M-1}\|D^{-1}\Delta u_i\|_2^2.
\]
The existing componentwise correction boxes imply $0\le R\le1$. The former
$R\le\sigma\le1$ epigraph is therefore eliminated exactly, giving
\[
 \min\ \bar\rho+\epsilon R(\Delta U),\qquad
 \sum_i\xi_i\le J_s^\star+\epsilon_s.
\]
The endpoint and CLF cones remain. At 76 holds, this removes one scalar and a
154-dimensional regularization cone: the primary has 630 variables and the
secondary 631. Explicit states retain the sparse dynamics structure. Inputs
are penalized about their linearization values, not about zero.

A small C++ MEX adapter calls Clarabel 0.11.1 with one thread. Equalities,
linear rows and finite bounds, and Lorentz cones are translated without
changing their feasible sets. Each call constructs a fresh sparse solver;
primal/dual warm starts or factorization reuse are not claimed. Full convergence
uses the configured desired accuracy. Reduced-accuracy convergence uses the
existing $10^{-5}$ numerical feasibility tolerance for residuals and objective
gap; flags 1 and 2 distinguish these outcomes. This numerical error is separate
from the analytic lexicographic tie bounds. Infeasibility certificates and
unfinished or unresolved iterates do not supply controls.

Strictly forcing reduced tolerances to equal the desired $10^{-8}$ precision
caused numerical stalls in otherwise useful boundary problems. Changing KKT
regularization globally did not repair all cases. The final adapter retains
the library's default linear-system regularization and records solver status,
iterations and residuals explicitly.

\subsection{Resolved-improvement stopping}
An accurate two-stage iterate with positive CLF slack at its numerical
first-input boundary continues to refine while either the predicted or actual
slack reduction exceeds $\epsilon\max(1,V)$. Below that existing objective
resolution, optional polishing ends. This is a local numerical stopping rule,
not a global nonlinear optimality proof. It does not stop merely because PCBF
slack is zero or the first pair is accurate.

The previous contraction on poor model agreement, next-frame trust-scale
adaptation, and bounded primary-infeasibility expansion remain. The new
accurate-step within-frame growth is not retained. The best accurate complete
pair from the current frame survives later failed trials; this is optimization
bookkeeping, not a previous-plan execution policy. One controller and one
nominal cost-to-go CLF remain.

\section{Paired original worst-frame experiment}
Every row uses the same recorded state, target epoch and previous trajectory
from the original 15 m/s braking-lead frame at 5.90 s. Four calls are made per
implementation; the median of the last three is reported. Source paths are
asserted. An earlier path-shadowed comparison that accidentally ran the final
source under every label is retained as an invalid diagnostic and excluded.
\begin{center}\small\begin{tabular}{lrrr}\toprule
Implementation & Seconds & Rounds & Numerical calls\\\midrule
Original & 2.948780 & 8 & 16\\
Analytic primary completion & 2.450923 & 8 & 10\\
Native quadratic conic solver & 0.433825 & 8 & 16\\
Added trust growth (rejected) & 0.315747 & 5 & 12\\
Retained stopping rule; no added growth & 0.427211 & 8 & 16\\
\bottomrule\end{tabular}\end{center}
The native solver selects numerically different future inputs within the same
convex problem. Subsequent nonlinear anchors need not be bitwise identical;
some strict zero-anchor certificates available in the MATLAB sequence are
unavailable along the native sequence. Algebraic equivalence does not imply
identical closed-loop trajectories.

All sixteen frozen programs return full-accuracy convergence. After restoring
the eliminated epigraph, their maximum original-constraint residual is
$1.2293\times10^{-11}$ and maximum objective difference from the stored MATLAB
points is $1.1797\times10^{-6}$. This and the independent known-QP/SOC tests
check conversion correctness separately from physical control performance.

\section{Final fourteen-scenario campaign}
Each prescribed encounter would collide under constant-speed given-path
cruising. Speeds are 8 and 15 m/s, prefix lengths eight and sixteen holds,
with a three-second completion horizon and a 50 ms hold. Target tangential
acceleration and sideslip are constant, states are observed exactly, and road
constraints are absent. The optimization-node buffer is 0.10 m; the physical
collision criterion is strictly positive rectangle clearance.

Recovery requires all five transverse errors to stay within
$[0.1\ \mathrm m,1^\circ,0.1\ \mathrm{m/s},0.05\ \mathrm{m/s},
0.01\ \mathrm{rad/s}]$ for one second after the prescribed interaction.
An observation cap is not a success criterion. ODE45 replay uses 31 samples
per hold, relative tolerance $10^{-11}$ and absolute tolerance $10^{-12}$.
Independent Python rectangle geometry and recovery-dwell checks agree.

\begin{center}\small\begin{tabular}{rlrrr}\toprule
Speed & Scenario & Minimum gap (m) & Recovery (s) & Peak call (s)\\\midrule
8 & headOn & 0.239939 & 12.05 & 1.263946\\
8 & acceleratingHeadOn & 0.053534 & 12.90 & 0.343604\\
8 & brakingLead & 0.097127 & 15.40 & 0.392661\\
8 & crossing & 0.086508 & 13.65 & 0.360544\\
8 & turningCrossing & 0.051406 & 14.10 & 0.342060\\
8 & curvedHeadOn & 0.469780 & 13.15 & 0.363576\\
8 & curvedCrossing & 0.049297 & 17.05 & 0.387361\\
15 & headOn & 0.085520 & 10.40 & 0.372675\\
15 & acceleratingHeadOn & 0.079247 & 8.75 & 0.317363\\
15 & brakingLead & 0.095731 & 13.75 & 0.514561\\
15 & crossing & 0.088047 & 12.80 & 0.369218\\
15 & turningCrossing & 0.000000 & 12.70 & 0.384337\\
15 & curvedHeadOn & 0.056895 & 66.95 & 0.362592\\
15 & curvedCrossing & 0.074375 & 14.45 & 0.362655\\
\bottomrule\end{tabular}\end{center}
The final run contains 4,762 issued holds and 147,622 independently audited
geometry samples. It records 2,403 analytic primary completions, 14,267
numerical calls and 368 reduced-accuracy convergences. Every issued frame
completes the CLF stage; the largest accepted normalized prediction discrepancy
is 0.995272. There are 594 calls longer than 0.1 s.

The circular 15 m/s head-on fixture has another encounter near 56 s. Its
60 s observation is resumed from the saved state until post-second-encounter
recovery at 66.95 s. This must not be interpreted as the duration of its first
avoidance maneuver. Recovery speed was not imposed as an acceptance criterion;
for example, 8 m/s braking-lead recovery changes from 14.20 to 15.40 s.

\subsection{Largest running frame}
Excluding only the process's first controller call, the maximum occurs in
15 m/s braking-lead at 6.10 s: 0.514561 s, eight rounds and eighteen numerical
calls. Its measured composition is
\begin{center}\begin{tabular}{lr}\toprule
Component & Seconds\\\midrule
PCBF solves & 0.122741\\
CLF solves & 0.160292\\
Formulation excluding CLF & 0.112811\\
CLF construction & 0.095964\\
Prediction agreement & 0.015188\\
Initialization & 0.004072\\
Other & 0.003493\\
\bottomrule\end{tabular}\end{center}
Numerical calls now occupy approximately 55\% of this frame; repeated
formulation and CLF construction occupy another 41\%. The original maximum
of eight rounds and soft five-second search budget are unchanged. No hard
0.1 s deadline is established. Controller latency is measured but not applied
to the simulated vehicle's command timing, so these tests do not establish
safe execution with delayed commands.

\section{Safety findings and rejected variants}
The final 15 m/s turning-crossing contact occurs between 1.485 and
1.48833 s in the hold starting at 1.45 s. Its start/mid/end gaps are
0.40603, 0.11233 and 0.10014 m, and its maximum constrained-node pose discrepancy
is only 0.1403 mm. The numerical node constraints can therefore hold while a
time between them contains overlap. The 0.10 m node buffer is not a
continuous-interval certificate. Faster solves and better linearization alone
cannot provide the missing interval constraint.

The full trust-growth trial finished all fourteen cases and recovered nominal
behavior, but only twelve were collision-free. Its additional curved-crossing
contact had positive start/mid/end gaps of 0.09998, 0.09746 and 0.34085 m, with
1.590 mm node-pose error. A further trial requiring positive predicted CLF
reduction before growth introduced another turning-crossing collision. It too
was rejected. Native-only controls on curved crossing achieve 0.074375 m
clearance; the native-only turning case still collides, matching the retained
version. Thus the stopping rule is not the source of that remaining contact.

The final full campaign uses the retained source throughout. Trial campaigns,
source snapshots, numerical failures and the invalid benchmark are preserved
separately. No scene-dependent controller selection or buffer adjustment is
used to improve a pass count. The remaining requirements are all-case interval
safety and an acceptable maximum runtime; neither is claimed achieved.

\section{Validation, provenance and reproduction}
All 268 selected MATLAB tests pass, covering two-stage priority, analytic
primary completion, positive-slack behavior, nominal recovery, model agreement,
flow initialization, geometry/model kernels, configuration, source budget and
six native-program interface cases. This is not a full-repository test claim.
Factory Code Analyzer settings report pre-existing sparse assembly and
unused-variable advisories. The new adapter/build/tests have no MATLAB Code
Analyzer messages.

The existing Clarabel.cpp dependency is under \texttt{solver/clarabel};
\path{scripts/buildPredictiveConicSolver.m} reads its CMake feature settings
and links the tested 0.11.1 library. Vendor sources and generated binaries are
excluded from the project commit. On this host, MATLAB R2026a's post-link
manifest validator also rejects a minimal MEX example. The build handles only
that diagnostic and requires the newly generated binary to load and solve an
independent known QP. Other compilation or smoke failures remain errors.

\path{scripts/runTimedFlowCampaign.m} and the unchanged nonlinear replay
harness generate the traces. Compact results, tests, source/artifact hashes
and trial summaries accompany this report; full MAT/JSON traces and exact
benchmark snapshots are preserved as export copies in the external archive.
The workstation was not isolated as a real-time target. All reported timings
are observations, not execution-time bounds.
\end{document}
